La capa de ``juego propiamente dicho'' es, a nivel código, prácticamente idéntica al TP x86, porque ya era agnóstica a
la ISA (punteros a estructuras, enteros, dibujos de pantalla). Los archivos \texttt{riscv/game.cpp} y
\texttt{riscv/sched.cpp} conservan la lógica de \texttt{game\_move}, \texttt{game\_take}, \texttt{game\_get\_id},
\texttt{sched\_spawn} y \texttt{sched\_kill}.

Las dos diferencias relevantes respecto al informe x86:

\begin{itemize}
  \item \textbf{\texttt{game\_move}} ya no hace \texttt{copy\_page} ni \texttt{mmu\_mapPage} manuales: todo eso está
  encapsulado en \texttt{mmu\_task\_move} (sección~\ref{sec:sv39}). El cuerpo queda enfocado exclusivamente en la lógica
  del juego (cálculo del delta, wrap vertical, espejado para el jugador B, chequeo de base, puntaje, redibujo).

  \item \textbf{\texttt{sched\_spawn}} no registra ningún descriptor de TSS en la GDT (no hay). El \emph{caller} en
  \texttt{kernel.cpp::spawn\_miner} se ocupa de:
  \begin{enumerate}
    \item reservar el slot y marcarlo como vivo (responsabilidad de \texttt{sched\_spawn});
    \item resetear el \texttt{TrapFrame} del slot (\texttt{init\_task\_if\_needed}); y
    \item llamar a \texttt{mmu\_task\_init(slot, player, x, y)}, que construye (o actualiza) la page-table y vuelca el
    binario del minero en la celda de spawn.
  \end{enumerate}
\end{itemize}

\begin{codesnippet}
\begin{verbatim}
 static void spawn_miner(uint32_t player_id) {
     int32_t slot = sched_spawn(player_id);
     if (slot < 0) return;
     const uint32_t u = (uint32_t)slot;
     g_tasks[u].initialized = 0;
     init_task_if_needed(u);
     const uint32_t x = sched_miners_state[u].pos_x;
     const uint32_t y = sched_miners_state[u].pos_y;
     mmu_task_init(u, player_id, x, y);
 }
\end{verbatim}
\end{codesnippet}

El \emph{cooldown} de spawn (\texttt{g\_spawn\_cooldown\_a/b}, 6 ticks) es un detalle UX agregado al port (evita
spawnear 20 mineros apretando una tecla sostenida); no tiene contraparte en el TP3 original.
